\documentclass{article}

\usepackage{tabularx}
\usepackage{booktabs}

\title{Reflection and Traceability Report on \progname}

\author{\authname}

\date{}

\input{../Comments.text}
\input{../Common.text}

\begin{document}

\maketitle

This document summarizes how feedback from Revision 0 was addressed in Revision 1
and reflects on key engineering decisions and outcomes. In addition to documenting
traceability of changes, it records lessons learned across project design,
verification, management, and professional practice outcomes.

\section{Changes in Response to Feedback}

Across Revision 1, feedback from the instructor, peer reviewers. The mapping of feedback to actions is summarized in the GitHub issue tracker, we will provide links to the related issues.

\subsection{SRS and Hazard Analysis}
\begin{itemize}
  \item \textbf{Instructor feedback:} The feedback was addressed by the instructor and peer reviewers. \href{https://github.com/zhangwenchili/CAS741-EelbrainPipeline/issues/23}{\#23}
  \item \textbf{Peer reviewer feedback:} \href{https://github.com/zhangwenchili/CAS741-EelbrainPipeline/issues/4}{\#4}, \href{https://github.com/zhangwenchili/CAS741-EelbrainPipeline/issues/5}{\#5}, \href{https://github.com/zhangwenchili/CAS741-EelbrainPipeline/issues/6}{\#6}, \href{https://github.com/zhangwenchili/CAS741-EelbrainPipeline/issues/7}{\#7}, \href{https://github.com/zhangwenchili/CAS741-EelbrainPipeline/issues/8}{\#8}
\end{itemize}

\subsection{Design and Design Documentation}

\begin{itemize}
    \item \textbf{Instructor feedback:} The feedback was addressed by the instructor and peer reviewers. \href{https://github.com/zhangwenchili/CAS741-EelbrainPipeline/issues/25}{\#25}
    \item \textbf{Peer reviewer feedback:} \href{https://github.com/zhangwenchili/CAS741-EelbrainPipeline/issues/14}{\#14}, \href{https://github.com/zhangwenchili/CAS741-EelbrainPipeline/issues/15}{\#15}, \href{https://github.com/zhangwenchili/CAS741-EelbrainPipeline/issues/16}{\#16}, \href{https://github.com/zhangwenchili/CAS741-EelbrainPipeline/issues/17}{\#17}, \href{https://github.com/zhangwenchili/CAS741-EelbrainPipeline/issues/19}{\#19}, \href{https://github.com/zhangwenchili/CAS741-EelbrainPipeline/issues/20}{\#20}, \href{https://github.com/zhangwenchili/CAS741-EelbrainPipeline/issues/21}{\#21}, \href{https://github.com/zhangwenchili/CAS741-EelbrainPipeline/issues/22}{\#22}
  \end{itemize}

\subsection{VnV Plan and Report}

\begin{itemize}
    \item \textbf{Instructor feedback:} The feedback was addressed by the instructor and peer reviewers. \href{https://github.com/zhangwenchili/CAS741-EelbrainPipeline/issues/24}{\#24}
    \item \textbf{Peer reviewer feedback:} \href{https://github.com/zhangwenchili/CAS741-EelbrainPipeline/issues/9}{\#9}, \href{https://github.com/zhangwenchili/CAS741-EelbrainPipeline/issues/10}{\#10}, \href{https://github.com/zhangwenchili/CAS741-EelbrainPipeline/issues/11}{\#11}, \href{https://github.com/zhangwenchili/CAS741-EelbrainPipeline/issues/12}{\#12}, \href{https://github.com/zhangwenchili/CAS741-EelbrainPipeline/issues/13}{\#13}
  \end{itemize}

\section{Extras}

The extras set out intially (user manual and code walkthrough) are cancelled by the instructor.

\section{Design Iteration (LO11 (PrototypeIterate))}

The design evolved from a functionality-first prototype to a traceability-oriented
engineering artifact set. Initially, the focus was enabling end-to-end M/EEG
processing by wrapping MNE-Python workflows. After reviews, the design was iterated
to make assumptions explicit, isolate likely changes behind module secrets, and
enforce requirement-to-test traceability.

The client need was rapid, repeatable analysis on BIDS datasets while keeping
developer effort manageable for future changes (e.g., additional BIDS entities and
preprocessing operations). This motivated the final decomposition where the pipeline
interface (M1) orchestrates behavior-hiding modules (M2--M6), and external
computation is delegated to MNE-Python (M7). The resulting design better supports
both usability for researchers and maintainability for developers.

\section{Design Decisions (LO12)}

\textbf{Limitations.}
Project scope and timeline limited full re-validation of external numerical
algorithms. Therefore, the team treated MNE-Python as a trusted dependency and
verified interface contracts and workflow integration rather than re-proving
algorithm correctness.

\textbf{Assumptions.}
The SRS assumptions (dipole approximation, linear relationship, and statistical noise
model) enabled concise instance models for inverse estimation. These assumptions
guided requirements and constrained what the software promises to users.

\textbf{Constraints.}
The BIDS input constraint and optional FreeSurfer MRI support strongly shaped the
architecture. BIDS compliance is checked/inferred early in the pipeline, and
module responsibilities were separated to localize changes in path inference,
preprocessing defaults, and covariance/source workflows.

\section{Economic Considerations (LO23)}

\progname{} is best positioned as an open-source research software component rather
than a commercial product. The likely user base includes labs and students working
with M/EEG data and already using MNE-Python/BIDS workflows.

Instead of direct sales, value comes from adoption, reproducibility, and reduced
analysis setup time. To attract users, the project should prioritize clear
documentation, example workflows, and stable automated tests.

\section{Reflection on Project Management (LO24)}

This section reflects on process and tooling choices used to deliver and revise the
project artifacts.

\subsection{How Does Your Project Management Compare to Your Development Plan}

The project broadly followed the intended plan: iterative document development,
review-driven revisions, and issue-based tracking in the repository. Planned tooling
(GitHub for version control/issues, checklist-guided reviews, and CI-oriented
testing strategy) was used as expected. The main variance from the original plan was
that more effort was needed for cross-document consistency and traceability updates
after feedback.

\subsection{What Went Well?}

\begin{itemize}
  \item Checklist-based reviews gave concrete, actionable feedback.
  \item Shared document templates improved consistency across SRS, MG, MIS, and VnV.
  \item Early use of traceability matrices helped identify missing links quickly.
  \item External supervisor/domain feedback improved realism of requirements and tests.
\end{itemize}

\subsection{What Went Wrong?}

\begin{itemize}
  \item Initial drafts optimized for progress speed, resulting in later rework for
  traceability and measurable verification criteria.
  \item Some requirement-test links were initially implicit rather than explicit.
  \item Synchronizing updates across multiple documents was time-consuming and
  increased the risk of inconsistency.
\end{itemize}

\subsection{What Would you Do Differently Next Time?}

\begin{itemize}
  \item Define traceability artifacts (requirement IDs, test IDs, module links) from
  the first draft instead of retrofitting them later.
  \item Schedule formal mini-walkthroughs earlier for high-risk sections
  (requirements clarity and module boundaries).
  \item Conduct periodic consistency checks across SRS/MG/MIS/VnV after each major edit.
\end{itemize}

\section{Ethics and Equity (LO18)}

The team emphasized equitable access by designing a scriptable and documented
pipeline that can be used on common operating systems (Linux/macOS/Windows) with
open tooling. Requirements and tests were written to prioritize reproducibility and
transparent error reporting, reducing hidden barriers for new users.

From a research ethics perspective, the project does not create or alter participant
data collection protocols; it operates on datasets prepared under existing domain
standards. The documentation-first approach also supports fair onboarding for users
with different prior experience levels.



\end{document}